\section{Flat scales carrying many singular centers}
\label{sec:flat-scales}

The first reduction uses only stationary monotonicity and Hausdorff
measure theory. If the singular set has dimension greater than $m-1$,
we shall find a singular point with a planar tangent cone and retain,
at the same scales, a nonzero measure of nearby points having exactly
the same density. Retaining this measure is essential: a planar
tangent cone at a single singular point gives no contradiction.

For $p\in U$ and $r>0$, write
\[
 \eta_{p,r}(X)=\frac{X-p}{r},\qquad
 V_{p,r}=(\eta_{p,r})_\#V.
\]
Thus the weight of $V_{p,r}$ is $r^{-m}(\eta_{p,r})_\#\mu$.
After translating $p$ to the origin we abbreviate $V_{p,r}$ to $V_r$.
We use $\mathbf B_r$ for ambient balls and, later, $B_r$ for balls in
an $m$-dimensional base plane. The stationary density is
\[
 \Theta_V(p)=\lim_{r\downarrow0}
       \frac{\mu(\mathbf B_r(p))}{\omega_m r^m}.
\]
It exists at every support point, is at least one, and is locally
bounded and upper semicontinuous. Moreover, under locally mass-bounded
convergence of stationary integral varifolds,
\begin{equation}
 \label{eq:density-upper-limit}
 V_i\rightharpoonup V_\infty,\quad p_i\longrightarrow p
 \quad\Longrightarrow\quad
 \Theta_{V_\infty}(p)\ge\limsup_i\Theta_{V_i}(p_i).
\end{equation}
These facts follow from monotonicity and integral compactness
\cite{Allard1972,Simon1983}. For example, to obtain
\eqref{eq:density-upper-limit}, bound each density by its mass ratio
on a fixed ball, pass to the limit at radii whose boundary has zero
limiting mass, and then let the radius decrease to zero.

We shall also use two consequences of these facts. First, supports
converge locally under this varifold convergence. The lower density
bound prevents support points from persisting where the limit has
zero mass; conversely, any neighborhood of a limiting support point
has positive limiting mass and therefore meets the approximating
supports. Second, $\spt\mu$ has locally finite $\Hh^m$ measure.
Indeed the lower density bound gives
$\mu(\mathbf B_r(p))\ge\omega_m r^m$ whenever the ball is interior
and $p\in\spt\mu$. A disjoint-ball covering argument on a compact
subset of $U$ bounds its Hausdorff measure by a dimensional constant
times the mass in a slightly larger compact set. No local
Hausdorff-measure assumption on the support is needed.

\begin{proposition}[Flat scales and a measure of centers]
\label{prop:flat-scales}
Suppose that $V$ is a stationary integral $m$-varifold in an open
subset of $\R^{m+n}$ and that
\[
 m-1<s<\dim_{\Hh}\Sing(V).
\]
After a translation there exist a fixed point $0\in\Sing(V)$, an
integer $Q\ge1$, an $m$-plane $\Pi$ through zero, radii $L_i\downarrow0$,
and positive Radon measures $\sigma_i$ such that
\begin{equation}
 \label{eq:flat-targets}
 \Theta_V(0)=Q,\qquad V_{L_i}\rightharpoonup Q|\Pi|
 \quad\hbox{locally in }\R^{m+n}.
\end{equation}
Each $\sigma_i$ is supported in $\overline{\mathbf B}_1$ on singular
points of $V_{L_i}$ having density exactly $Q$. There are constants
$c,C>0$, independent of $i$, for which
\begin{equation}
 \label{eq:center-measure-bounds}
 \sigma_i(\mathbf B_u(X))\le C u^s
 \quad(X\in\R^{m+n},\ u>0),\qquad
 c\le\sigma_i(\overline{\mathbf B}_1)\le C.
\end{equation}
\end{proposition}

\begin{proof}
The regular set is relatively open in $\spt\mu$, so the singular set
is Borel. The compact-subset theorem for Hausdorff measure
\cite[Theorem~8.13]{Mattila1995} gives a compact subset of
$\Sing(V)$ with finite positive $\Hh^s$ measure.
We recall the two measure-theoretic refinements that we need
\cite{Federer1969,Mattila1995}. On a set of finite positive
$\Hh^s$ measure, the upper $s$-density is positive and finite almost
everywhere. By restricting to a positive-measure subset where its
upper bound is uniform, and then taking a compact subset, we obtain
a compact set $K$ such that
\begin{equation}
 \label{eq:initial-frostman}
 0<\Hh^s(K)<\infty,\qquad
 \Hh^s(K\cap\mathbf B_r(X))\le C r^s.
\end{equation}
An upper bound initially obtained for centers in $K$ extends to
arbitrary centers by doubling the radius. Enlarging $C$ handles the
remaining large radii. Lusin's theorem permits a further compact
restriction of positive measure on which $\Theta_V$ is continuous.
The upper bound in \eqref{eq:initial-frostman} is preserved.

For $\Hh^s$-almost every $p\in K$, positive upper density supplies
radii $r_i\downarrow0$ and a constant $c_p>0$ with
\[
 \Hh^s(K\cap\mathbf B_{r_i}(p))\ge c_p r_i^s.
\]
Pass to a subsequence on which $V_{p,r_i}$ converges to a stationary
integral tangent cone $C_p$. Define measures on the closed unit ball by
\[
 \lambda_i=r_i^{-s}(\eta_{p,r_i})_\#
       \big(\Hh^s\!\llcorner(K\cap\overline{\mathbf B}_{r_i}(p))\big).
\]
They have uniformly bounded mass, mass bounded away from zero, and
the upper ball bound in \eqref{eq:initial-frostman}. A further
subsequence converges to a nonzero measure $\lambda$ with the same
upper ball bound. Its support lies in $\spt C_p$.

Set $q_p=\Theta_V(p)$; at this point $q_p$ is not assumed to be an
integer. If $z\in\spt\lambda$, there are points $p_i\in K$ with
$(p_i-p)/r_i\to z$. Continuity on $K$ and
\eqref{eq:density-upper-limit} imply
\[
 \Theta_{C_p}(z)\ge q_p.
\]
On the other hand, the density at infinity of $C_p$ about $z$ equals
$q_p$: sandwich $\mathbf B_R(z)$ between the balls about zero of
radii $R-|z|$ and $R+|z|$, and use conicality about zero. Monotonicity
about $z$ consequently has equality at every radius.

We spell out the geometric consequence of this equality. The radial
displacement $X-z$ lies in the tangent plane of $C_p$ almost
everywhere. Testing stationarity with $\varphi(X)(X-z)$ gives
\[
 \int\big(m\varphi+(X-z)\cdot\nabla\varphi\big)
       \,\dd\|C_p\|=0.
\]
Differentiating dilated scalar integrals shows that the weight is
$m$-homogeneous about $z$. Since a rectifiable weight determines its
tangent plane almost everywhere, the varifold is conical about $z$
as well. Composing dilations about $z$ and about zero now gives
translation invariance along the whole line $\R z$.

The support of $\lambda$ spans at least $m$ linearly independent
directions. Otherwise it lies in a linear space of dimension at most
$m-1$; covering its bounded support by $O(\varepsilon^{-(m-1)})$
balls of radius $\varepsilon$ would give
$\lambda(\R^{m+n})\le C\varepsilon^{s-m+1}\to0$.
Choose $m$ independent such directions and let $P$ be their span.
The cone is invariant under every translation in $P$. If its support
contained a point outside $P$, subtracting its $P$-component would
give $w\in P^\perp\setminus\{0\}$ in the support. Conicality and
translation invariance would then place every disk
$tw+(P\cap\mathbf B_1)$, $1\le t\le2$, in the support. These are
disjoint disks of equal positive $m$-dimensional measure, all in a
bounded region, contradicting local finiteness of
$\Hh^m(\spt C_p)$. Thus its support is precisely $P$.
Tangential stationarity on this plane gives constant multiplicity,
and integrality gives
\[
 C_p=q_p|P|,\qquad q_p\in\{1,2,\ldots\}.
\]

We have proved that $\Theta_V$ is integer-valued at
$\Hh^s$-almost every point of $K$. One of its countably many integer
level sets therefore has positive $\Hh^s$ measure. Choose a compact
positive-measure subset $K_Q$ of that level set. Repeating the
preceding argument with $K_Q$, choose a fixed point $p\in K_Q$ and
radii $L_i\downarrow0$ on which its normalized $s$-mass is bounded
below. The resulting tangent cone is a $Q$-plane. Translate $p$ to
zero, denote that plane by $\Pi$, and set
\[
 \sigma_i=L_i^{-s}(\eta_{p,L_i})_\#
       \big(\Hh^s\!\llcorner
          (K_Q\cap\overline{\mathbf B}_{L_i}(p))\big).
\]
Compactness of $K_Q$ ensures that the support consists entirely of
points from that integer-density level set. Dilation preserves both
density and singularity. The asserted convergence and bounds follow.
\end{proof}
